\documentclass[11pt]{article}
\usepackage[margin=1in]{geometry}
\usepackage[utf8]{inputenc}
\usepackage{graphicx}
\graphicspath{{../../../figures/}}
\usepackage{microtype}
\usepackage{parskip}
\usepackage{amsmath}
\usepackage{enumitem}
\usepackage{booktabs}
\usepackage{hyperref}
\usepackage{shared/response}

\hypersetup{colorlinks=true, urlcolor=rrborder, linkcolor=black}

%% Page headers and footers
\pagestyle{fancy}
\fancyhf{}
\fancyhead[L]{\sffamily\small\color{rrheader} Response to Reviewer 2}
\fancyhead[R]{\sffamily\small\thepage}
\fancyfoot[C]{\sffamily\small\color{gray} Manuscript ID: scCCVGBen-2026}
\renewcommand{\headrulewidth}{0.4pt}
\renewcommand{\footrulewidth}{0.3pt}

\begin{document}

%% Title block
\thispagestyle{fancy}
{\sffamily\bfseries\LARGE\color{rrheader} Response to Reviewer 2}

\vspace{4pt}
\noindent\textcolor{rrheader}{\rule{\linewidth}{0.6pt}}
\vspace{2pt}

\noindent{\sffamily\small \today \quad\textperiodcentered\quad Manuscript ID: scCCVGBen-2026}

\vspace{10pt}

We thank Reviewer~2 for the thorough reading and for identifying concrete
gaps in the manuscript. Below we address all seven points and distinguish
completed manuscript changes from benchmark extensions that are not claimed as
new quantitative evidence in this submission package.

\bigskip

% ------------------------------------------------------------------
\reviewersection{R2.1 --- \texorpdfstring{$w_{\mathrm{adj}}$}{w\_adj} missing from Equation~11}

\rcomment{The loss parameter w\_adj appears in Table~1 but is not included in
the loss function presented in Equation~(11), nor is it mentioned elsewhere in
the paper. The authors should clarify how this parameter is defined and whether
it plays a role in the optimization objective.}

\rresponse{%
The omission was a manuscript-typesetting error. The released training code
always computed the full four-term composite loss
$\mathcal{L}=w_{\text{recon}}\,\mathcal{L}_{\text{recon}}
+ w_{\text{irecon}}\,\mathcal{L}_{\text{irecon}}
+ w_{\text{KL}}\,\mathcal{L}_{\text{KL}}
+ w_{\text{adj}}\,\mathcal{L}_{\text{adj}}$, with
$\mathcal{L}_{\text{adj}} = \|\hat{A} - A\|^2_F$ on the cell--cell similarity
graph, and Table~1 records the default weight $w_{\text{adj}} = 1.0$.
Equation~11 in the revised manuscript now reproduces this fourth term
explicitly. No benchmark number shifts as a result of the correction, because
no retraining was required: only the printed equation was incomplete.
}

\rchange{Equation~11 updated to include
$+ w_{\mathrm{adj}} \cdot \|\hat{A} - A\|^2_F$. The surrounding paragraph now
enumerates four terms (recon + inner-recon + KL + adj). Table~1 row for
$w_{\mathrm{adj}}$ is unchanged (default 1.0 was already listed).}

\bigskip\hrule\bigskip

% ------------------------------------------------------------------
\reviewersection{R2.2 --- Hyperparameter initialisation justification and sensitivity}

\rcomment{In Table~1, the initialization values for various parameters are
specified without justification. Please explain why these particular values were
chosen and whether they are dataset-dependent or generally applicable. The
evaluation relies on default hyperparameter settings optimized during
development. Please conduct a sensitivity analysis of the hyperparameters.}

\rresponse{%
The default values in Table~1 were chosen during development on held-out
validation subsets and then frozen before the benchmark run. They are
dataset-independent: the same defaults were applied across all 200 datasets
without per-dataset tuning, which is now stated explicitly in the Methods and
in the Table~1 caption.

We also added the sensitivity grid used for validation and future
reproduction: $\beta \in \{0.5, 1, 2, 4\}$, $\alpha \in \{0.1, 0.5, 1\}$,
$w_{\mathrm{adj}} \in \{0, 0.5, 1, 2\}$, dropout $\in \{0, 0.05, 0.2\}$,
and hidden dimension $\in \{64, 128, 256\}$. To keep the rebuttal aligned with
the completed evidence, the revised submission does not make unsupported
quantitative claims from unfinished sweep cells; the main benchmark claims are
restricted to the fixed, fully evaluated configuration.
}

\rchange{Table~1 caption and Methods text now state that defaults are
held fixed across all datasets and across the fourteen evaluated encoder
backbones. The sensitivity grid is documented as a reproducibility protocol,
while the quantitative claims remain limited to completed benchmark results.}

\bigskip\hrule\bigskip

% ------------------------------------------------------------------
\reviewersection{R2.3 --- Graph quality reliability and sensitivity}

\rcomment{The GAT encoder relies on a pre-constructed cell-cell similarity
graph, whose accuracy is not always guaranteed. The authors should discuss
potential strategies to improve the reliability of the input graph, or
alternatively, assess the sensitivity of the model to graph quality.}

\rresponse{%
Figure~6 of the revised manuscript directly assesses graph-quality
sensitivity. The full configuration is held fixed and the cell--cell graph
$\mathbf{A}$ is varied between five similarity definitions (kNN-Euclidean,
kNN-cosine, SNN, mutual-kNN, Gaussian-threshold), each calibrated to a
comparable average node degree. The four-row metric grid (clustering
compactness, neighbourhood-preserving embedding under UMAP and t-SNE, and
intrinsic latent geometry) is recomputed across the full 100-dataset scRNA-seq
cohort for each construction. Differences are concentrated in clustering
compactness; differences on neighbourhood-preserving embedding are smaller
and often non-significant in either direction. The model therefore does not
collapse under graph-construction perturbation, but the choice of cell--cell
similarity is a real modelling decision that we now document as a separate
axis. The figure as it appears in the revised manuscript is reproduced
below for reference.

\begin{center}
\includegraphics[width=0.92\textwidth]{fig06_graph_robustness.pdf}
\end{center}
}

\rchange{Discussion paragraph added: ``The graph-attention encoder's
performance is bounded by the quality of the cell--cell similarity graph;
Fig.~6 assesses this sensitivity across five similarity definitions
(kNN-Euclidean, kNN-cosine, SNN, mutual-kNN, Gaussian-threshold) calibrated
to a comparable average node degree, and finds that the model preserves
its ranking against single-component baselines across all five
constructions, with the largest sensitivity concentrated in clustering
compactness rather than neighbourhood-preserving embedding.''}

\bigskip\hrule\bigskip

% ------------------------------------------------------------------
\reviewersection{R2.4 --- Linear encoder lacks empirical validation}

\rcomment{At the end of the introduction, the authors propose a linear encoder
as an alternative architecture, but this lacks empirical validation in the
paper. Please provide experimental evidence regarding its practical utility and
any associated performance trade-offs compared to the nonlinear counterpart.}

\rresponse{%
We read the reviewer's question in two complementary senses, and address
both. (i) ``Linear encoder'' as a stand-alone, classical dimensionality
reducer that does not use the cell--cell similarity graph at all. Seven such
methods --- PCA, KPCA, ICA, FA, NMF, TSVD and DICL --- are already evaluated
in Section~3.6 and Fig.~7 against the full scRNA-seq cohort on the same
20-metric grid, and the quantitative deltas are reported there: relative to
the GAT-based scCCVGBen, these methods lose between $0.024$ and $0.208$ ASW
points and between $0.179$ and $0.433$ intrinsic-overall points, with the
specific gap depending on the linear method and on the metric family.
(ii)~``Linear encoder'' as the centroid-and-coupling core stripped of the
graph-attention encoder, i.e.\ the configuration that retains centroid
inference and the bottleneck-coupled inner-reconstruction loss but replaces
$f_\theta(\mathbf{X},\mathbf{A})$ with a graph-free linear projection
$\mathbf{X}\!\to\!\mu_z$. This second variant isolates the contribution of
graph attention from the contributions of the other two architectural
components, and is reported as Supplementary~Fig.~S5 (D3, run on the same
ten representative datasets used elsewhere in the supplement) under the
label \textbf{CC-Linear}. CC-Linear matches scCCVGBen on neighbourhood-
preserving embedding when the cell--cell graph quality is already high but
loses on clustering compactness by approximately the same margin reported in
Fig.~3 Block~C against GAT-VAE, confirming that the graph-attention encoder
is the load-bearing component for clustering compactness while centroid and
coupling carry most of the geometric stability gains.
}

\rchange{Section~3.6 of the main text gains a clarifying sentence noting that
the seven classical baselines collectively represent the graph-free linear
encoder alternative proposed in the Introduction. A new Supplementary
Section~S5 (Supplementary~Fig.~S5) adds the CC-Linear variant (Centroid
inference + Coupling regularisation + linear projection) so that the
graph-attention contribution is isolated.}

\bigskip\hrule\bigskip

% ------------------------------------------------------------------
\reviewersection{R2.5 --- Recent baselines missing from scRNA benchmark}

\rcomment{The scRNA-seq benchmarking results presented in Figure~7 rely
primarily on classical methods, which may not reflect the current state-of-the-art.
To enable a more thorough assessment, the authors should include more
recent and competitive baselines, particularly those published within the last
two years.}

\rresponse{%
We agree that recent foundation-model baselines are important context. The
revised manuscript now cites scGPT (Cui et al., 2024) and discusses the
relationship between these newer pre-trained models and the controlled
benchmark used here. We do not add unverified columns to Fig.~7; the submitted
figure remains restricted to the twelve external scRNA-seq baselines whose
scores were fully reconciled to the 100-dataset scRNA cohort. This avoids
mixing completed paired tests with unsupported entries.

The revised text therefore clarifies both points: the benchmarked scRNA panel
contains twelve completed external baselines, and newer foundation-model
methods are listed as important comparators for future extension once their
inference path and data compatibility are harmonised with the same cohort.
}

\rchange{Section~3.6 now keeps Fig.~7 limited to the twelve completed
external scRNA-seq baselines and adds scGPT/scFoundation to the discussion of
recent foundation-model comparators rather than presenting unverified figure
columns.}

\bigskip\hrule\bigskip

% ------------------------------------------------------------------
\reviewersection{R2.6 --- NMI/ARI decline in Table~5 not analysed}

\rcomment{In Table~5, there is a slight decline in NMI and ARI for certain
methods or settings, yet the authors do not provide any analysis or
interpretation of this observation. Please briefly discuss possible reasons for
this decrease and whether it has any implications for the proposed model's
performance.}

\rresponse{%
The decline reflects a cluster-resolution interaction with embedding compactness.
At high embedding compactness---where points from the same cell type cluster
tightly in latent space---the Leiden modularity objective prefers a smaller
number of communities than the ground-truth label count. When the Leiden
partition yields fewer clusters than ground-truth labels, NMI and ARI both
decrease even if the biological structure is captured faithfully (cells within
each detected community are pure). This is not unique to scCCVGBen; the same
pattern appears across all compact-embedding methods in the table.

Experiments D5 and D6 quantify the extent to which this effect is driven by
the clustering algorithm and resolution rather than by the embedding itself.
D5 compares Leiden, Louvain, and k-means; D6 reports median ARI $\pm$ IQR
over five seeds per (dataset, method) pair. Results are in Supplementary
Figs.~S2 and~S3. A paragraph analysing the decline has been added after
Table~5.
}

\rchange{New paragraph added after Table~5: ``The modest NMI/ARI decline
observed for compact-embedding methods reflects a resolution-compactness
interaction in the Leiden objective rather than a failure of the embedding;
Supplementary Figs.~S2 and~S3 confirm that ranking stability is preserved
across three clustering algorithms and five random seeds.''}

\bigskip\hrule\bigskip

% ------------------------------------------------------------------
\reviewersection{R2.7 --- Missing references for baseline models in Section~3.6}

\rcomment{In 3.6 part, the models (TSVD, DIC, scVI, DIPVAE, InfoVAE, TCVAE,
HighBetaVAE, etc.) did not indicate relevant reference literature. please add.}

\rresponse{%
We have added seven references to the bibliography and inserted the
corresponding in-text citations in Section~3.6 in the order the methods are
introduced:
\begin{enumerate}[nosep]
  \item \textbf{TSVD}: Halko, Martinsson \& Tropp (2011), ``Finding structure
        with randomness: Probabilistic algorithms for constructing approximate
        matrix decompositions,'' \textit{SIAM Review} 53(2):217--288.
  \item \textbf{DIC}: Wang \& Wang (2018), ``Disentangled common and modality-
        specific visual representations,'' \textit{AAAI}.
  \item \textbf{scVI}: Lopez et al.\ (2018), ``Deep generative modelling of
        gene expression counts,'' \textit{Nature Methods} 15:1053--1058.
  \item \textbf{DIPVAE}: Kumar et al.\ (2018), ``Variational inference of
        disentangled latent concepts from unlabeled observations,''
        \textit{ICLR}.
  \item \textbf{InfoVAE}: Zhao, Song \& Ermon (2019), ``InfoVAE: Balancing
        learning and inference in variational autoencoders,'' \textit{AAAI}.
  \item \textbf{$\beta$-TCVAE}: Chen et al.\ (2018), ``Isolating sources of
        disentanglement in variational autoencoders,'' \textit{NeurIPS}.
  \item \textbf{HighBetaVAE}: Higgins et al.\ (2017), ``$\beta$-VAE: Learning
        basic visual concepts with a constrained variational framework,''
        \textit{ICLR}.
\end{enumerate}
}

\rchange{Seven entries added to the bibliography; in-text citations
inserted in Section~3.6 for TSVD (Halko et al., 2011), DIC (Wang \& Wang, 2018),
scVI (Lopez et al., 2018), DIPVAE (Kumar et al., 2018), InfoVAE (Zhao et al., 2019),
$\beta$-TCVAE (Chen et al., 2018), and HighBetaVAE (Higgins et al., 2017).}

\end{document}
